\chapter{Related Work}

\section{Linear-time alternatives to softmax attention}

The quadratic time and memory cost of softmax
attention~\cite{vaswani2017attention} has motivated a sustained
line of work on linear-time alternatives. A first wave
approximates the softmax operator with low-rank or kernel feature
maps, yielding sub-quadratic transformers such as
Performer~\cite{choromanski2021performer},
Linformer~\cite{wang2020linformer}, and
Reformer~\cite{kitaev2020reformer}. These approaches preserve the
explicit pairwise attention abstraction while reducing its
computational footprint.

A second, more architecturally divergent line replaces softmax
attention with a recurrent state. Linear
Attention~\cite{katharopoulos2020transformers} reformulated
softmax attention as a causal recurrence with bounded state, and
subsequent work refined the recurrence structure through gating
(GLA~\cite{yang2024gla}), rank-one Householder updates
(DeltaNet~\cite{schlag2021deltanet,yang2024deltanet}), retention
masks (RetNet~\cite{sun2023retnet}), and matrix-valued cells
(xLSTM~\cite{beck2024xlstm}). State-space models (SSMs) reach the
same operating point from a continuous-time root: the S4
family~\cite{gu2022s4} introduced structured diagonal-plus-low-rank
parameterisations, Mamba~\cite{gu2023mamba} added input-selective
dynamics, and Mamba-2~\cite{dao2024transformers} unified SSMs and
linear attention through structured state-space duality;
Mamba-3~\cite{lahoti2026mamba3} continues this lineage. Long
convolutional alternatives such as Hyena~\cite{poli2023hyena}
share the same linear-complexity envelope through implicit
kernels rather than an explicit recurrence.

The RWKV family occupies a parallel branch built directly on a
recurrent core. RWKV-4~\cite{peng2023rwkv} introduced the
time-mix and channel-mix decomposition, RWKV-5 and RWKV-6 (Eagle
and Finch~\cite{peng2024eaglefinch}) added matrix-valued states
and dynamic per-channel decay, and RWKV-7~\cite{peng2025rwkv7}
extends the transition class to a diagonal-plus-low-rank form.
The thesis instantiates three architectures from this collective
body of work as concrete probes for its mechanism analysis:
causal RWKV-6 as a per-channel data-dependent decay backbone,
Mamba-2 as a selective continuous-time SSM with input-modulated
step size, and Linear Attention as a pure accumulator without
intrinsic decay. The three share a recurrent core and linear-time
token mixing while differing in the structure of their transition
operator and in the strength of their content-adaptive routing.
The remainder of this chapter positions the mechanisms studied
in the thesis against the deficits these structural differences
introduce.

\section{Bidirectional adaptation strategies for causal linear RNNs}

Causal linear-time recurrences propagate state in one direction,
which suits streaming and autoregressive workloads but discards
future context that is available in offline settings such as
document encoding, full-utterance speech recognition, and image
classification. Three structurally distinct strategies have
emerged for adapting an existing causal architecture to
bidirectional input access. At the mechanism level,
LION~\cite{afzal2025linear} replaces the causal scan kernel with
a parallel $T \times T$ attention computation under a symmetric
decay mask, retaining the time-mix and channel-mix abstraction
while changing the sequence-access pattern from causal to
bidirectional. At the operator level,
Vision-RWKV~\cite{duan2024visionrwkv} introduces a bidirectional
WKV scan that runs the causal RWKV recurrence forwards and
backwards and combines the resulting hidden states; Vision
Mamba~\cite{zhu2024visionmamba} applies the same template to
selective state-space models, pairing forward and backward Mamba
blocks within each layer. At the ordering level, no new operator
is introduced: the input is rearranged so that successive layers
see the sequence in different traversal orders, as in
VisualRWKV's per-layer scan alternation~\cite{hou2024visualrwkv},
VMamba's four-route Cross-Scan Module~\cite{liu2024vmamba},
LocalMamba's windowed selective scan with dynamic per-layer
pattern selection~\cite{huang2024localmamba}, and ZigMa's
zero-parameter zigzag pattern~\cite{hu2024zigma}. The three
strategies trade off compute, memory, and information-flow
geometry: the mechanism-level form pays a quadratic memory
footprint in exchange for unrestricted future access at every
layer; the operator-level form retains a linear-time envelope but
doubles the serial compute per layer; the ordering-level form
costs only a tensor reshape, but bidirectional information flow
accumulates across the depth stack rather than within any single
layer.

The thesis adopts the mechanism-level strategy and instantiates
LION as a parallel adaptation across all three linear-time
backbones. This choice is motivated by the structural property
that the bidirectional form inherits the causal mechanism
vocabulary directly, and is supported by an empirical comparison
against the operator-level alternative reported in the
experimental chapter.

\section{Cross-modality short-range temporal mixing}

A documented failure mode of vanilla linear-time sequence mixers
is the absence of an explicit local mixing operator: each token
sees only its own embedding plus the running recurrent state, and
that state cannot encode short-range structure as efficiently as
an explicit convolution can. The cost of this deficit is most
visible in the original Vision Mamba
paper~\cite{zhu2024visionmamba}, which transferred causal
selective state-space models to image classification but reported
scaling difficulties beyond modest parameter counts.
Vision-RWKV~\cite{duan2024visionrwkv}, published in parallel,
attached a fixed quad-directional shift (Q-Shift) to each token
before the recurrent core and scaled stably past 30~M parameters
on the same benchmark. The structural difference is a single
per-block operator that gives every token explicit access to its
four spatial neighbours.

\begin{table}[t]
\centering
\small
\caption[Cross-modality local-mixing operators]{Local-mixing
operators introduced on top of linear-time recurrent backbones.
$x \in \mathbb{R}^{T \times d}$ denotes the per-block input,
$V \in \mathbb{R}^{T \times d}$ the value tensor of a linear
attention layer, $\mathrm{DWConv}_{k, d}$ a depth-wise 1D
convolution of kernel size $k$ and dilation $d$, and
$\mathrm{DWSep2D}_{k}$ a depth-wise separable 2D convolution
of kernel size $k$.}
\label{tab:local-mixing-cross-domain}
\setlength{\tabcolsep}{4pt}
\renewcommand{\arraystretch}{1.08}
\begin{tabularx}{\linewidth}{@{}p{2.25cm} p{2.45cm} p{1.45cm} X@{}}
\toprule
\textbf{Mechanism} & \textbf{Reference} & \textbf{Backbone} & \textbf{Operator} \\
\midrule
Q-Shift &
Vision-RWKV~\cite{duan2024visionrwkv} &
RWKV-4/5 &
Channels split into four groups; each group shifted by $\pm 1$ in
one of $\{$up, down, left, right$\}$; concatenated before the
time-mix block \\
\addlinespace
Depth-wise $V$-conv &
LiT~\cite{wang2025lit} &
Linear DiT &
$V \leftarrow \mathrm{DWConv}_{k}(V)$ inserted on the value path
of linear attention \\
\addlinespace
Cross-Scan Module &
VMamba~\cite{liu2024vmamba} &
Mamba &
Image flattened along four canonical routes; SSM scan applied to
each; the four outputs merged \\
\addlinespace
Windowed scan &
LocalMamba~\cite{huang2024localmamba} &
Mamba &
SSM applied independently within non-overlapping image windows;
per-layer dynamic pattern selection \\
\addlinespace
Local conv prefix &
U-Mamba~\cite{ma2024umamba}, MambaIR~\cite{guo2024mambair} &
Mamba &
$x \leftarrow \mathrm{DWConv}_{k}(x)$ prepended to every block \\
\addlinespace
QuadScan &
U-RWKV~\cite{ye2025urwkv} &
RWKV &
Four-route raster scan with direction selection; generalises the
fixed Q-Shift \\
\addlinespace
2D DWSep ConvShift &
AudioRWKV~\cite{xiong2025audiorwkv} &
RWKV-7 &
$x_{\mathrm{local}} = \mathrm{DWSep2D}_{k}(x)$ added as a
residual feeding the decay, key, and value paths \\
\addlinespace
Multi-dilation DWConv &
Non-Attention LLM~\cite{kiruluta2025breaking} &
Linear attention &
$\sum_{d \in \{1, 2, 4, 8\}} \alpha_{d} \,
\mathrm{DWConv}_{k, d}(x)$ on the input side \\
\bottomrule
\end{tabularx}
\end{table}

The local-mixing remediation has been independently rediscovered
across the linear-time literature, and the operators differ in
syntactic form rather than function.
Table~\ref{tab:local-mixing-cross-domain} summarises the eight
principal instances. Within the RWKV lineage, Q-Shift carries
through subsequent vision and vision-language
models~\cite{fei2024diffusionrwkv,hou2024visualrwkv,gu2024rwkvclip},
generalises into a direction-adaptive QuadScan in U-RWKV's
medical segmentation backbone, and reappears alongside a
learnable two-dimensional depth-wise separable ConvShift on
RWKV-7 in AudioRWKV. Within the SSM lineage, VMamba's Cross-Scan
Module delivers an analogous four-route access pattern from the
SSM side; LocalMamba restricts the scan to local windows; and
U-Mamba together with MambaIR prepend an explicit local
convolution to every block. On the diffusion side, LiT attaches
a depth-wise convolution to the value path of a distilled linear
DiT. In autoregressive language modelling, the Non-Attention
LLM of Kiruluta et al. proposes a parallel bank of dilated
depth-wise convolutions on the input side with dilations $\{1, 2,
4, 8\}$, the most direct precedent for the input-side multi-scale
convolution introduced in this thesis. The dilation set covers
receptive fields appropriate for token-level structure in
language and, after the acoustic subsampling that precedes the
encoder, for the phoneme-to-syllable scale in speech.

The pattern spans two-dimensional vision in both the RWKV and SSM
lineages, two-dimensional spectrogram audio on a
diagonal-plus-low-rank RWKV-7 backbone, and one-dimensional
autoregressive language on a linear-attention backbone, before
reappearing in this thesis on one-dimensional speech. The
recurring observation is that linear-complexity sequence mixers,
regardless of the algebraic structure of their transition
operator, leave a short-range mixing deficit that an explicit
depth-wise-convolution-class operator restores. Audio
Mamba~\cite{erol2024audiomamba} provides the cleanest
counterpoint: a self-attention-free bidirectional state-space
model for audio classification that does not introduce explicit
local mixing, leaving the deficit unaddressed in that work. This
cross-domain robustness is the empirical foundation for the
input-side multi-scale convolution mechanism studied in the
proposed-solution chapter.

\section{Associative memory and recall}

The capacity to retrieve a value associated with a previously
seen key is one of the operations that softmax attention performs
naturally and that linear-time recurrent architectures struggle
to match at scale. The structural reason is direct: softmax
attention compares every query against every key in parallel,
while a linear-time recurrent model compresses the history into
a bounded state. When several stored keys are similar, their
values overlap in the state and become difficult to disentangle
at readout. The Zoology study of Arora et
al.~\cite{arora2023zoology} formalised this gap with the
multi-query associative recall (MQAR) benchmark, in which a
sequence of $(k, v)$ pairs is followed by query keys whose values
must be retrieved; sub-quadratic models degrade markedly as the
number of stored pairs approaches the state dimension. The same
gap has been corroborated through fixed-state copy
failures~\cite{jelassi2024repeat}, through in-context retrieval
as the core bottleneck remaining after chain-of-thought
prompting~\cite{wen2024rnns}, and through in-context-learning
studies showing that Mamba matches transformers on most tasks
but lags on retrieval-heavy
ones~\cite{park2024mamba,akyurek2024incontext}. Mechanistically,
the induction-head circuit~\cite{olsson2022induction} provides
an account of how softmax attention realises recall through
pairwise key matching, clarifying what an architectural
alternative must reproduce.

Three families of architectural responses have followed. The
first preserves softmax attention in part of the model: Based of
Arora et al.~\cite{arora2024based} interleaves a Taylor expansion
of softmax with linear attention to balance recall and
throughput, while GoldFinch~\cite{goldstein2024goldfinch} pairs
an RWKV-6 backbone with a transformer post-processor under heavy
KV-cache compression. The second modifies the recurrent update
of a linear-time backbone to actively cancel interference.
DeltaNet, introduced by Schlag et al.~\cite{schlag2021deltanet}
as a fast-weight programmer, applies a rank-one Householder
update $(I - \beta_t k_t k_t^\top)$ that erases the component of
the state aligned with the current key before writing the new
value, interpretable as one online gradient step on the squared
associative reconstruction error. Yang et
al.~\cite{yang2024deltanet} parallelise this update over sequence
length and combine it with input-dependent gating to obtain
Gated DeltaNet, integrating the delta-rule erasure with the
selective gating mechanism characteristic of the SSM lineage.
The Compact Recurrent Transformer of Mucllari et
al.~\cite{mucllari2025crt} pursues a related approach with a
Cayley-orthogonal transition. The third family operates within
softmax attention itself, sharpening retrieval without altering
the quadratic substrate. LUCID of Duvvuri et
al.~\cite{duvvuri2026lucid} replaces the standard
$\mathrm{softmax}(QK^\top/\sqrt{d}) V$ output with
$\mathrm{softmax}(QK^\top/\sqrt{d}) \cdot P^{-1} V$, where
$P = M \circ \exp(KK^\top)$ is a preconditioner constructed from
key-key similarities; the construction can be derived as a
delta-rule update in the reproducing kernel Hilbert space induced
by the exponential kernel, decorrelating values along directions
of high key correlation before they enter the attention output.
The mechanism preserves the $O(N^2 d)$ complexity of softmax
attention and targets retrieval quality on long-context tasks
rather than computational cost.

These three response families address recall capacity but do not
all engage the formal expressivity of the recurrent state.
DeltaNet's rank-one update moves the architecture along a
complexity-theoretic axis, lifting it out of the strictly
diagonal class to which the other linear recurrences belong. The
next section turns to the formal side of this question,
presenting the circuit-complexity results that delimit what
diagonal and diagonal-plus-low-rank recurrences can and cannot
compute.

